\title{Controlling Text-to-Image Diffusion by Orthogonal Finetuning}

\begin{abstract}
Large-scale text-to-image diffusion models demonstrate remarkable proficiency in synthesizing photorealistic images from natural language prompts. Effectively steering or modulating these advanced models for a variety of downstream applications remains a prominent open question. In response, we propose a systematic finetuning technique termed Orthogonal Finetuning~(OFT), designed to tailor text-to-image diffusion models for specialized tasks. Distinct from prevailing approaches, OFT is theoretically guaranteed to maintain hyperspherical energy, which encapsulates the relationships between neuron pairs on the unit hypersphere. Our investigations reveal that retaining this property is essential for upholding the semantic generation capability inherent to text-to-image diffusion frameworks. To further enhance the robustness of finetuning, we introduce Constrained Orthogonal Finetuning (COFT), incorporating an extra constraint on the hypersphere's radius. Our study focuses on two significant finetuning scenarios: subject-driven generation, where the objective is to render images depicting a specific subject in various contexts based on a limited set of subject images and a text prompt; and controllable generation, which seeks to incorporate auxiliary control signals into the generation process. Experimental evaluations demonstrate that our OFT approach surpasses current state-of-the-art methods in both the quality of generated images and training convergence rate.
\end{abstract}

\section{Introduction}

State-of-the-art text-to-image diffusion models~\cite{saharia2022photorealistic,ramesh2022hierarchical,rombach2022high} have set new benchmarks in producing high-fidelity images conditioned on textual descriptions. However, despite their strong baseline performance, relying solely on text input often leaves room for ambiguity and restricts precise and nuanced manipulation of generated outputs. To mitigate these limitations, this work focuses on two central classes of text-to-image generation tasks:

\begin{itemize}[leftmargin=*,nosep]

    \item \textbf{Subject-driven generation}~\cite{ruiz2023dreambooth}: This problem involves, given a small number of images of a particular subject, generating novel images of that same subject in varied settings, leveraging both the visual examples and the provided textual prompt.
    \item \textbf{Controllable generation}~\cite{zhang2023adding,mou2023t2i}: Here, alongside a text prompt, an explicit control signal (\emph{e.g.}, canny edges, segmentation maps) is provided to direct the generative process, aiming to closely align the output image with the given control.
\end{itemize}

At their core, both tasks address the challenge of how to adapt text-to-image diffusion models through finetuning while maintaining the generative capabilities gained during pretraining. We distill the essential requirements for a robust finetuning approach into two primary aspects: (1) \emph{training efficiency}, characterized by minimizing the number of parameters that require updating and reducing the total number of training epochs; and (2) \emph{generalizability preservation}, referring to the retention of the model's capacity to generate high-quality and diverse outputs. 

To achieve these aims, standard finetuning practices either adjust the model’s existing neuron weights using a modest learning rate (\emph{e.g.}, \cite{ruiz2023dreambooth}) or introduce a small auxiliary module consisting of re-parameterized neuron weights (\emph{e.g.}, \cite{hulora2022,zhang2023adding}). Although these strategies are straightforward, neither method can robustly ensure that the generative performance attained during pretraining is conserved. Moreover, a clear framework for developing effective finetuning strategies and for determining optimal hyperparameters—such as the appropriate number of epochs—remains elusive.

A fundamental challenge stems from the absence of a direct metric for assessing the degree to which a finetuned model preserves the pretrained generative abilities. Contemporary finetuning methodologies often operate under the implicit premise that a smaller Euclidean distance between the weights of the finetuned model and those of the pretrained model signals superior preservation of pretraining. For this reason, finetuning generally employs extremely low learning rates. While this heuristic occasionally yields satisfactory results, we contend that the Euclidean distance by itself does not adequately reflect the extent of semantic retention. Therefore, adopting a more structurally-informed metric to evaluate the differences between the finetuned and pretrained models could substantially enhance both the retention of generative performance and the overall stability of the finetuning process.

Drawing on the empirical findings that hyperspherical similarity effectively encapsulates semantic relationships~\cite{liu2017deep,liu2018decoupled,chen2020angular}, our approach utilizes hyperspherical energy~\cite{liu2018learning} as a metric to represent the pairwise relational dynamics among neurons within a layer. The concept of hyperspherical energy entails aggregating the hyperspherical similarity (such as cosine similarity) for every pair of neurons residing in the same layer, thereby measuring how uniformly the neurons are distributed on the unit hypersphere~\cite{liu2021learning}. We put forward the conjecture that an optimal finetuned model should maintain a minimal discrepancy in hyperspherical energy when compared to its pretrained counterpart. A straightforward approach might be to introduce a regularization term to stabilize hyperspherical energy throughout finetuning, but this does not ensure effective minimization of the energy gap. 

To address this, we exploit a key property of hyperspherical energy: when all neurons undergo the same orthogonal transformation, pairwise hyperspherical similarity remains invariant. Based on this property, we introduce Orthogonal Finetuning~(OFT), a method designed to adapt large-scale text-to-image diffusion models for downstream applications without altering their hyperspherical energy. The main strategy involves learning an orthogonal transformation, shared across the entire layer, that preserves the angles between neuron vectors. This process can be interpreted as recalibrating the standard coordinate basis of the layer’s neurons. Recognizing that maintaining a small Euclidean distance between the finetuned and original pretrained models is beneficial for preserving performance characteristics, we further develop a variant known as Constrained Orthogonal Finetuning~(COFT). In COFT, the finetuned parameters are restricted to remain within a hyperspherical shell of constant radius, centered on the positions of the pretrained neurons.

The rationale behind employing orthogonal transformations for neuron finetuning is partly inspired by the principles of the 2D Fourier transform, which decomposes an image into magnitude and phase spectra. Among these, the phase spectrum encapsulates the angular relationship between the input and the basis functions and is responsible for retaining most of the semantic information. Illustratively, reconstructing an image using its phase spectrum alongside an arbitrary magnitude spectrum can often preserve the essential semantics of the original. This observation implies that altering neuron orientations plays a pivotal role in making meaningful, semantic modifications to generated images, aligning with the objectives of subject-driven and controllable generation tasks. Nevertheless, unconstrained modifications to neuron directions can severely compromise the generative performance established during pretraining. To regulate this flexibility, we propose maintaining the angles between every pair of neurons, premised on the belief that these inter-neuronal angles are vital for the neural network’s internalized knowledge. This leads to a natural design choice: applying layer-wide orthogonal transformations to the neurons in each layer, thereby conserving hyperspherical energy.

Additionally, our approach takes cues from orthogonal over-parameterized training~\cite{liu2021orthogonal}, which develops classification neural networks by introducing orthogonal transformations to the initial random weights. Randomly initialized networks typically exhibit low expected hyperspherical energy, and the work of \cite{liu2021orthogonal} aims to preserve this property during training, as reduced hyperspherical energy has been linked to improved generalization in classification~\cite{liu2018learning,lin2020regularizing}. They demonstrate that orthogonal transformations are sufficiently expressive to enable the training of highly generalizable classifiers. By contrast, our methodology targets the finetuning of text-to-image diffusion models to achieve enhanced controllability and generative capacity for downstream tasks. We stress two primary distinctions between OFT and the method proposed in \cite{liu2021orthogonal}. Firstly, while \cite{liu2021orthogonal} endeavors to minimize hyperspherical energy, OFT is designed to preserve the hyperspherical energy characteristic of the pretrained model, thereby safeguarding the model’s inherent structure during finetuning—minimizing hyperspherical energy in this context could disrupt original semantic representations. Secondly, OFT explicitly aims to minimize divergence from the pretrained parameters, resulting in a more constrained finetuning variant, whereas \cite{liu2021orthogonal} operates without such restrictions. Finetuning inherently involves negotiating the tension between adaptability and structural integrity, and we contend that OFT strikes a judicious balance between these competing demands. Our principal contributions are enumerated below:

\begin{itemize}[leftmargin=*,nosep]

    \item We introduce a new finetuning technique called Orthogonal Finetuning, which is designed to enhance the controllability of text-to-image diffusion models. Additionally, to bolster training stability, we present a constrained version that restricts the model's angular divergence from its original pretrained state.
    \item Unlike current finetuning approaches, OFT adjusts the model parameters while mathematically maintaining the hyperspherical energy—an attribute we empirically recognize as vital for preserving the generative semantics learned by the pretrained model.
    \item Our approach is applied to two scenarios: subject-driven image synthesis and controllable generation. Through thorough experiments, we show that OFT delivers superior results relative to existing methods, exhibiting improved image quality, faster convergence, and greater finetuning robustness. Notably, OFT also demonstrates increased sample efficiency, achieving effective convergence with significantly fewer training images and iterations.
    \item In the context of controllable generation, we propose a novel metric for assessing control consistency. This metric involves inferring the control input from the synthesized image and then measuring its alignment with the original control directive. The use of this metric further substantiates the enhanced controllability provided by OFT.
\end{itemize}

\section{Related Work}

\textbf{Text-to-image diffusion models}. The landscape of text-to-image synthesis has rapidly advanced~\cite{saharia2022photorealistic,ramesh2022hierarchical,rombach2022high,gu2022vector,nichol2022glide}, largely owing to breakthroughs in diffusion-based generative modeling~\cite{ho2020denoising,song2021score,song2021denoising,dhariwal2021diffusion} as well as joint vision-language representation learning~\cite{su2020vlbert,lu2019vilbert,radford2021learning,wang2021simvlm,li2022blip,li2023blip,alayrac2022flamingo,schuhmann2022laion}. GLIDE~\cite{nichol2022glide} and Imagen~\cite{saharia2022photorealistic} operate by training diffusion models directly on pixel-level data. While GLIDE integrates its text encoder training with a diffusion prior leveraging paired data, Imagen opts for a frozen, pretrained text encoder. Diffusion in latent space is explored in models such as Stable Diffusion~\cite{rombach2022high} and DALL-E2~\cite{ramesh2022hierarchical}, where the former utilizes VQ-GAN~\cite{esser2021taming} to construct a visual latent codebook and the latter employs CLIP~\cite{radford2021learning} for creating a joint representation space of text and images. Beyond diffusion architectures, approaches rooted in generative adversarial networks~\cite{reed2016generative,zhang2017stackgan,xu2018attngan,li2019controllable} and autoregressive modeling~\cite{ramesh2021zero,ding2021cogview,wu2022nuwa,yuscaling} are also actively explored. OFT is designed as a universally applicable finetuning mechanism suitable for any text-to-image diffusion model.

\textbf{Subject-driven generation}. Approaches that utilize user-supplied masks as extra conditioning to avoid subject alteration include~\cite{avrahami2022blended,nichol2022glide}. Various inversion-based strategies~\cite{choi2021ilvr,dhariwal2021diffusion,ramesh2022hierarchical,gal2022image} target background editing while maintaining the subject intact. Local and global modifications without explicit mask input are made feasible by~\cite{hertz2022prompt}. Yet, these methods often fail to retain crucial identity-specific features. Pivotal Tuning~\cite{roich2022pivotal} addresses identity retention by finetuning the generator near the inverted latent representation, incorporating additional regularization terms. A similar objective is pursued by~\cite{nitzan2022mystyle}, which builds a tailored generative facial prior based on images of a given subject. Although \cite{casanova2021instance} is able to produce various renditions of an instance, it can sometimes overlook distinguishing details. DreamBooth~\cite{ruiz2023dreambooth} adapts text-to-image diffusion models using a reconstruction loss and a class-based prior preservation loss, leveraging a limited number of subject-specific images and an assigned custom token. Our OFT method operates within the DreamBooth framework; however, rather than straightforward parameter finetuning at a low learning rate, OFT introduces parameter updates via orthogonal transformations.

\textbf{Controllable generation}. Image-to-image translation represents a paradigm of controllable generation, where most existing approaches leverage conditional generative adversarial networks~\cite{isola2017image,zhu2017toward,wang2018high,park2019semantic,choi2018stargan}. Additionally, diffusion models have emerged as competitive tools for this task~\cite{saharia2022palette,voynov2022sketch,wang2022pretraining}. Recently, ControlNet~\cite{zhang2023adding} introduced a mechanism to exert control over a pretrained diffusion model by finetuning it with supplementary control inputs, resulting in remarkable advancements in controllable generation. In parallel, T2I-Adapter~\cite{mou2023t2i} proposes a similar strategy, wherein a pretrained diffusion model undergoes finetuning to enhance controllability over the output images. In alignment with the settings described in \cite{zhang2023adding,mou2023t2i}, our method employs OFT to finetune pretrained diffusion models, consistently outperforming in controllability while utilizing reduced training data and fewer finetuned parameters. Importantly, OFT incurs no additional inference-time computation cost.

\textbf{Model finetuning}. Recently, the practice of finetuning large pretrained models for downstream applications has seen growing adoption~\cite{devlin2018bert,brown2020language,he2022masked}. Within the broader landscape of finetuning, adaptation techniques (\emph{e.g.}, \cite{hulora2022,houlsby2019parameter,pfeiffer2020adapterfusion}) are extensively explored in natural language processing research. Of these, LoRA~\cite{hulora2022} shares the closest conceptual relationship to OFT, relying on a low-rank decomposition of weight modifications throughout finetuning. OFT, on the other hand, implements layer-wide orthogonal transformations that multiplicatively adjust neuron weights. This approach mathematically maintains the pair-wise angles among neurons within each layer, offering improved finetuning stability.

\section{Orthogonal Finetuning}

\subsection{Why Does Orthogonal Transformation Make Sense?}

We begin by examining the rationale behind employing orthogonal transformations in finetuning text-to-image diffusion models. Specifically, we frame the inquiry in two parts: (1) what is the motivation for adjusting the orientation (\emph{i.e.}, angle) of neurons during finetuning, and (2) what justifies the selection of orthogonal transformations for modifying these angles?

Addressing the first question, we build on the empirical insights from \cite{liu2018decoupled,chen2020angular}, which indicate that differences in feature angles effectively represent semantic disparities. SphereNet~\cite{liu2017deep} demonstrates that constraining all neural network activations to reside on a unit hypersphere maintains comparable expressiveness and even enhances generalization performance. This finding suggests that the orientation of neuronal vectors encapsulates essential data representations. To further illustrate the significance of neuron orientation, we present a simple experiment in Figure~\ref{fig:toy_ae}, wherein a conventional convolutional autoencoder is trained on images of flowers. During training, the feature map is computed using the standard inner product between the convolutional filter $\bm{w}$ (with element output $z$) and the local input $\bm{x}$. During testing, we compare three different approaches to computing feature maps: (a) the training inner product, (b) preserving only magnitude, and (c) retaining only angular information. The outcomes, as visualized in Figure~\ref{fig:toy_ae}, reveal that utilizing solely angular neuron information enables nearly flawless reconstruction of the original images, whereas magnitudes alone provide no semantic recovery. Notably, cosine-based activation (c) is not used in training—the training employs the standard inner product exclusively. These findings underscore that the directional component of neuronal parameters is chiefly responsible for encoding image semantics. Therefore, to effect meaningful changes to image semantics, adjusting the directions of neurons through targeted finetuning is likely the most effective approach.

Turning to the second question, a straightforward method for finetuning neuron directions involves collectively rotating or reflecting all neurons within a layer, which inherently utilizes orthogonal transformations. Although alternative angular modifications—such as varying rotation angles per neuron—could provide added flexibility, our observations suggest that orthogonal transformations strike an optimal balance between adaptability and regularity. Additionally, \cite{liu2021orthogonal} provides evidence that orthogonal transformations possess ample capacity for training neural networks. To substantiate this point, we conduct an experiment evaluating the regularizing effect of orthogonal constraints, following the experimental setup of ControlNet~\cite{zhang2023adding}. Figure~\ref{fig:ortho} displays these results, showing that our default OFT yields consistently stable performance and precise control once training concludes (after 20 epochs). Conversely, when the orthogonality constraint is omitted, OFT fails to produce plausible images or exert control, confirming the central role of the orthogonality constraint in the success of OFT.

\subsection{General Framework}

In traditional finetuning, gradient descent is applied with a relatively low learning rate to update the parameters of a model (or select layers), ensuring only slight modifications are made from the pretrained state. This is achieved because a small learning rate effectively restricts the extent of change, and as a result, standard finetuning implicitly targets minimization of $\thickmuskip=2mu \medmuskip=2mu \| \bm{M} - \bm{M}^0\|$, where $\bm{M}$ stands for the adapted model parameters, and $\bm{M}^0$ represents the original pretrained values. While this implicit regularization is conceptually reasonable, it may remain excessively permissive, especially when applied to large models. To counteract this limitation, LoRA introduces a direct low-rank limitation on weight adaptations, formally specified as $\thickmuskip=2mu \medmuskip=2mu \text{rank}(\bm{M} - \bm{M}^0)=r'$, where $r'$ is selected to be a small integer. In contrast, OFT enforces an alternative restriction, maintaining pairwise neuron similarity by requiring $\thickmuskip=2mu \medmuskip=2mu \| \text{HE}(\bm{M})-\text{HE}(\bm{M}^0) \|=0$, with $\text{HE}(\cdot)$ denoting the hyperspherical energy of a matrix. 

To clarify, let us examine the case of a fully connected layer $\thickmuskip=2mu \medmuskip=2mu \bm{W}=\{\bm{w}_1,\cdots,\bm{w}_n\}\in\mathbb{R}^{d\times n}$, where each neuron $\bm{w}_i\in\mathbb{R}^{d}$ (and pretrained weights are given as $\bm{W}^0$). The output, $\thickmuskip=2mu \medmuskip=2mu \bm{z}\in\mathbb{R}^n$, for an input $\thickmuskip=2mu \medmuskip=2mu \bm{x}\in\mathbb{R}^d$, is then computed as $\thickmuskip=2mu \medmuskip=2mu \bm{z}=\bm{W}^\top\bm{x}$. The OFT technique can thus be understood as encouraging the hyperspherical energy difference between the finetuned and original weights to be minimized:

\begin{equation}\label{eq:oft_principle}
\footnotesize
\min_{\bm{W}} \left\| \text{HE}(\bm{W})-\text{HE}(\bm{W}^0)\right\|~~\Leftrightarrow~~\min_{\bm{W}} \bigg{\|} \sum_{i\neq j} \|\hat{\bm{w}}_i-\hat{\bm{w}}_j\|^{-1}-\sum_{i\neq j} \|\hat{\bm{w}}^0_i-\hat{\bm{w}}^0_j\|^{-1}\bigg{\|}
\end{equation}

Here, the normalized neuron is represented by $\thickmuskip=2mu \medmuskip=2mu \hat{\bm{w}}_i:=\bm{w}_i/\|\bm{w}_i\|$, and the hyperspherical energy for $\bm{W}$ is calculated as $\thickmuskip=2mu \medmuskip=2mu \text{HE}(\bm{W}):= \sum_{i\neq j} \|\hat{\bm{w}}_i-\hat{\bm{w}}_j\|^{-1}$. Notably, the optimal (minimum) value for Eq.~\eqref{eq:oft_principle} is zero, attainable when $\bm{W}$ and $\bm{W}^0$ differ only by an orthogonal transformation—meaning $\thickmuskip=2mu \medmuskip=2mu \bm{W}=\bm{R}\bm{W}^0$ where $\thickmuskip=2mu \medmuskip=2mu \bm{R}\in\mathbb{R}^{d\times d}$ is orthogonal. The nature of $\bm{R}$'s determinant ($1$ or $-1$) determines if this is a rotation or a reflection. This underpins the key concept in OFT: rather than altering individual neurons

\begin{equation}\label{eq:oft_general}
    \footnotesize
    \bm{z} = \bm{W}^\top \bm{x} = (\bm{R} \cdot \bm{W}^0)^\top \bm{x},~~~\text{s.t.}~\bm{R}^\top\bm{R} = \bm{R}\bm{R}^\top = \bm{I}
\end{equation}

In this formulation, $\bm{W}$ is the weight matrix after OFT finetuning, while $\bm{I}$ represents the identity matrix. An overview of OFT is presented in Figure~\ref{fig:block}. As with the zero initialization procedure utilized in LoRA, it is essential for OFT that the finetuning process commences exactly from the pretrained parameter set $\bm{W}^0$. This is ensured by initializing the orthogonal matrix $\bm{R}$ as an identity matrix so that the initial state of the finetuned model coincides with the pretrained parameters. The orthogonality constraint on $\bm{R}$ can be maintained through differentiable orthogonalization methods as explored in \cite{liu2021orthogonal,lezcano2019cheap}. We will later elaborate on strategies to enforce orthogonality while keeping computations tractable.

\subsection{Efficient Orthogonal Parameterization}\label{sect:eff_ortho}

Although standard differentiable orthogonalization techniques such as the Gram-Schmidt procedure exist, their computational overhead is significant for large-scale applications~\cite{liu2021orthogonal}. To improve scalability, we utilize the Cayley parameterization to obtain orthogonal matrices. More concretely, we define the orthogonal matrix via $\thickmuskip=2mu \medmuskip=2mu \bm{R} = (\bm{I} + \bm{Q})(I - \bm{Q})^{-1}$, where $\bm{Q}$ is a skew-symmetric matrix, i.e., $\thickmuskip=2mu \medmuskip=2mu \bm{Q} = -\bm{Q}^\top$. While this method efficiently enforces orthogonality, it only generates matrices with determinant $1$—restricting outputs to the special orthogonal group. In our experiments, we observe that this limitation does not adversely impact practical performance. Even so, direct parameterization of a full-sized orthogonal matrix $\bm{R}$ remains inefficient when $d$ is large. To overcome this, we suggest representing $\bm{R}$ as a block-diagonal matrix, partitioned into $r$ blocks, so that:

\begin{equation}
    \footnotesize
    \bm{R} = \text{diag}(\bm{R}_1, \bm{R}_2, \cdots, \bm{R}_r) = \begin{bmatrix}
    \bm{R}_1 \in \text{O}(\frac{d}{r}) &  &\\
    &\ddots & \\
    & & \bm{R}_r \in \text{O}(\frac{d}{r})
    \end{bmatrix} \in \text{O}(d)
\end{equation}

Here, $\text{O}(d)$ represents the orthogonal group of dimension $d$, and both $\thickmuskip=2mu \medmuskip=2mu \bm{R}\in\mathbb{R}^{d\times d}$ and $\thickmuskip=2mu \medmuskip=2mu \bm{R}_i\in\mathbb{R}^{d/r\times d/r},\forall i$ are orthogonal matrices. In the specific case where $\thickmuskip=2mu \medmuskip=2mu r=1$, the block-diagonal orthogonal matrix collapses to a fully unconstrained orthogonal matrix. The total number of free parameters for a $d \times d$ orthogonal matrix is $\thickmuskip=2mu \medmuskip=2mu d(d-1)/2$, which corresponds to a parameter complexity of $\mathcal{O}(d^2)$. In contrast, an orthogonal matrix with $r$ blocks contains $\thickmuskip=2mu \medmuskip=2mu d(d/r-1)/2$ parameters, giving a complexity of $\thickmuskip=2mu \medmuskip=2mu \mathcal{O}(d^2/r)$. It is possible to reduce the parameter count further by tying together the block matrices, i.e., enforcing $\thickmuskip=2mu \medmuskip=2mu\bm{R}_i=\bm{R}_j,\forall i\neq j$. This brings the parameter complexity down to $\mathcal{O}(d^2/r^2)$. Notwithstanding these efficiency enhancements, it is important to highlight that the resulting matrix $\bm{R}$ still satisfies the orthogonality property, thereby fully maintaining hyperspherical energy.

To compare the parameter efficiency of OFT and LoRA, recall that for LoRA with a chosen low-rank $r'$, the count of trainable parameters is $\thickmuskip=2mu \medmuskip=2mu r'(d+n)$. Assuming both $r$ and $r'$ are linear in the neuron dimension $d$, such as $\thickmuskip=2mu \medmuskip=2mu r=r'=\alpha d$ for some constant $\thickmuskip=2mu \medmuskip=2mu 0<\alpha\leq 1$, the resulting parameter complexity for LoRA is $\mathcal{O}(d^2+dn)$, whereas for OFT it reduces to $\mathcal{O}(d)$. To make this more tangible, consider a weight matrix of dimensions $\thickmuskip=2mu \medmuskip=2mu 128\times 128$. If we set $\thickmuskip=2mu \medmuskip=2mu r'=8$ in LoRA, this results in $2,048$ trainable parameters. For OFT with $\thickmuskip=2mu \medmuskip=2mu r=8$ (without block sharing), the number of trainable parameters is $960$.

\subsection{Constrained Orthogonal Finetuning}

We can impose stricter constraints on OFT by requiring that the adapted model stay in close proximity to the pretrained model’s parameter space. To this end, COFT uses the following transformation at inference:

\begin{equation}\label{eq:coft_general}
    \footnotesize
    \bm{z}=\bm{W}^\top\bm{x}=(\bm{R}\cdot\bm{W}^0)^\top\bm{x},~~~\text{s.t.}~\bm{R}^\top\bm{R}=\bm{R}\bm{R}^\top=\bm{I},~\norm{\bm{R}-\bm{I}}\leq \epsilon
\end{equation}

This expression involves two types of constraints: $\bm{R}$ must be orthogonal and its deviation from the identity matrix is upper-bounded by $\epsilon$. The orthogonality requirement can be enforced using the Cayley parameterization, as detailed in Section~\ref{sect:eff_ortho}. However, including the $\epsilon$-bounded deviation constraint with the Cayley representation presents a nontrivial challenge. To further analyze the Cayley transform, we apply the Neumann series to approximate $\thickmuskip=2mu \medmuskip=2mu \bm{R}=(\bm{I}+\bm{Q})(I-\bm{Q})^{-1}$ as $\thickmuskip=

As the series converges in the operator norm, it is possible to rewrite the constraint $\thickmuskip=2mu \medmuskip=2mu\|\bm{R}-\bm{I}\|\leq\epsilon$ within the Cayley transform. This yields an equivalent requirement on the skew-symmetric matrix: $\thickmuskip=2mu \medmuskip=2mu\|\bm{Q}-\bm{0}\|\leq\epsilon'$, where $\bm{0}$ designates a matrix of all zeros and $\epsilon'$ serves as a different tolerance parameter from $\epsilon$. The new constraint imposed on $\bm{Q}$ is straightforwardly managed through projected gradient descent. To ensure the orthogonal matrix $\bm{R}$ starts as the identity, we simply set $\bm{Q}$ to be an all-zero matrix at initialization. COFT incorporates two explicit regularization terms: one enforcing minimal hyperspherical energy alteration and another controlling the distance from the pretrained model. While most standard finetuning strategies implicitly assume the latter, COFT formalizes it as an explicit objective. Although OFT achieves strong performance, our empirical results indicate that COFT delivers enhanced finetuning stability, attributable to the explicit control over model deviation. Figure~\ref{fig:eps_coft} illustrates the influence of $\epsilon$ on COFT's behavior. The figure demonstrates that the hyperparameter $\epsilon$ modulates the adaptation capacity: increasing $\epsilon$ allows COFT to emulate OFT, while reducing $\epsilon$ results in the COFT model acting more like the original, unmodified pretrained text-to-image diffusion model.

\subsection{Re-scaled Orthogonal Finetuning}

We introduce a straightforward modification to OFT by enabling the learning of a separate scaling parameter for each neuron. The rationale is that neuron rescaling does not impact hyperspherical energy, as magnitude normalization compensates for scaling. Formally, the forward computation is: $\thickmuskip=2mu \medmuskip=2mu\bm{z}=(\bm{R}\bm{W}^0\bm{D})^\top\bm{x}$\footnote{\textbf{Errata}: The NeurIPS camera-ready submission mistakenly listed the forward computation as $\thickmuskip=2mu \medmuskip=2mu\bm{z}=(\bm{D}\bm{R}\bm{W}^0)^\top\bm{x}$. However, the code uses the correct formula, so the reported results in Appendix~\ref{app:rescaled} are valid.}, where $\thickmuskip=2mu \medmuskip=2mu \bm{D}=\text{diag}(s_1,\cdots,s_n)\in\mathbb{R}^{n\times n}$ is a diagonal matrix with strictly positive, learnable entries $s_1,\ldots,s_n$. In this extension, both the orthogonal matrix $\bm{R}$ and the scaling matrix $\bm{D}$ are learned, as opposed to the baseline OFT formulation in Eq.~\eqref{eq:oft_general}, where only $\bm{R}$ is trainable. This augmentation boosts OFT's representational flexibility at the cost of an insignificant increase in parameters. Although the principal experiments utilize the original OFT to isolate the benefits of orthogonal transformations, our experiments show that incorporating re-scaling (see Appendix~\ref{app:rescaled}) tends to yield superior outcomes.

\section{Intriguing Insights and Discussions}

\textbf{OFT applies broadly across architectures}. OFT is, by design, suitable for any neural network architecture. In the context of Transformers, where LoRA is commonly employed to adapt attention weights~\cite{hulora2022}, we adopt a similar strategy: for a fair assessment, OFT is exclusively applied to attention parameters in our Transformer experiments. Moreover, OFT’s applicability extends naturally to convolutional networks. Thanks to the block-diagonal organization of $\bm{R}$, it affords interpretations in convolutional settings that are less straightforward for LoRA. If the number of blocks matches the number of input channels, then each block acts solely on its designated channel, analogous to tuning depth-wise convolutional filters~\cite{chollet2017xception}. Alternatively, employing shared blocks means that a single orthogonal transformation is applied across all channels. Additional results on convolutional layer finetuning using OFT are presented in Appendix~\ref{app:convolution}.

\textbf{Connection to LoRA}. LoRA stabilizes the pretrained weight matrix by incorporating a low-rank addition, thereby mitigating substantial shifts in its information. In comparison, OFT constrains the transformation applied to the pretrained weights to remain orthogonal (and thus full-rank), safeguarding the integrity of information learned during pretraining. The forward computation in OFT can be expressed as $\thickmuskip=2mu \medmuskip=2mu \bm{z}=(\bm{R}\bm{W}^0)^\top\bm{x}=(\bm{W}^0+(\bm{R}-\bm{I})\bm{W}^0)^\top\bm{x}$, where $\thickmuskip=2mu \medmuskip=2mu (\bm{R}-\bm{I})\bm{W}^0$ functions analogously to the low-rank update in LoRA. Assuming $\bm{W}^0$ has full rank, OFT also achieves a low-rank update provided that $\thickmuskip=2mu \medmuskip=2mu \bm{R}-\bm{I}$ is low-rank. Both approaches employ a hyperparameter to manage trainable parameter efficiency: LoRA uses a rank parameter $r'$, while OFT utilizes a diagonal block parameter $r$. Notably, the parameter-efficiency strategies differ fundamentally—LoRA leverages low-rankness, reducing trainable parameter count via low-rank structures; OFT, conversely, capitalizes on structural sparsity through block-diagonal orthogonal matrices.

\textbf{Why OFT converges faster}. Evidence from Figure~\ref{fig:toy_ae} suggests that the most impactful way to alter model semantics is by rotating neuron directions—a process that OFT explicitly facilitates. Furthermore, OFT effectively recasts finetuning as moving neurons on a smooth hypersphere manifold, thereby providing a more tractable optimization landscape. This property is also corroborated empirically by the findings in \cite{liu2021orthogonal}.

\textbf{Why not minimize hyperspherical energy}. Our approach diverges from \cite{liu2021orthogonal} by not targeting the minimization of hyperspherical energy. In classification tasks, it is preferable for neurons to avoid redundancy. Enforcing minimum hyperspherical energy uniformly distributes neurons across the hypersphere, a property that, while useful in some settings, does not serve finetuning objectives and can jeopardize valuable information from pretraining.

\textbf{Balancing flexibility and regularization in finetuning}. Our findings indicate a fundamental trade-off between flexibility and regularity during finetuning. While conventional finetuning methods offer maximal flexibility, they are prone to instability and can often result in model collapse. Notably, OFT, despite its straightforwardness, achieves an effective compromise by maintaining the pairwise angles between neurons. Additionally, the block-diagonal parameterization acts as a stronger regularizer for the orthogonal matrix.

\textbf{No extra inference computation}. In contrast to ControlNet, the OFT framework does not impose any extra computational burden during inference on the adapted model. During inference, one can directly integrate the learned orthogonal matrix $\bm{R}$ into the pretrained weights $\bm{W}^0$ to form the new weights $\thickmuskip=2mu \medmuskip=2mu \bm{W}=\bm{R}\bm{W}^0$, thereby maintaining the inference latency of the original pretrained model.

\section{Experiments and Results}

\textbf{Overall experimental setup}. Stable Diffusion v1.5~\cite{rombach2022high} serves as the pretrained backbone for text-to-image synthesis in our experiments. For equitable evaluation, sample images are randomly selected from each approach. Subject-driven tasks largely adhere to the DreamBooth~\cite{ruiz2023dreambooth} protocol, whereas controllable generation procedures follow ControlNet~\cite{zhang2023adding} and T2I-Adapter~\cite{mou2023t2i}. To directly compare against LoRA, both OFT and COFT are exclusively applied to those layers where LoRA operates. Detailed information and further empirical results are documented in Appendix~\ref{app:settings}.

\subsection{Subject-driven Generation}

\textbf{Experimental details}. DreamBooth~\cite{ruiz2023dreambooth} and LoRA~\cite{hulora2022} are employed as comparative baselines. All techniques utilize the identical loss formulation prescribed by DreamBooth. The implementations for DreamBooth and LoRA remain consistent with their respective original descriptions and utilize their recommended hyperparameters. Additional analyses and results can be found in Appendix~\ref{app:settings},\ref{app:coft_oft},\ref{app:more_qualitative},\ref{app:failure}.

\textbf{Finetuning stability and convergence}. We begin by assessing the stability and convergence rate during finetuning for DreamBooth, LoRA, OFT, and COFT, with the results summarized in Figure~\ref{fig:teaser} and Figure~\ref{fig:db_convergence}. Both COFT and OFT display stable behavior throughout the finetuning process. After 400 iterations, effective subject control is achieved with DreamBooth and OFT variants, whereas LoRA does not successfully maintain subject identity. At 2000 iterations, DreamBooth starts to produce degraded outputs, while LoRA fails to generate the yellow shirt and erroneously produces yellow fur instead. In contrast, OFT and COFT continue to deliver stable and reliable control over image synthesis. These outcomes provide strong evidence for the rapid and robust convergence characteristics of the OFT framework in subject-driven image generation tasks. Importantly, the performance of OFT and COFT shows little dependence on the specific number of finetuning iterations, reducing the need for manual adjustment of iteration counts per subject. This allows for a straightforward choice of a higher iteration number for both OFT and COFT, without the requirement for meticulous tuning. For COFT, selecting an appropriate $\epsilon$ parameter simplifies the configuration of both the learning rate and iteration count.

\textbf{Quantitative comparison}. In line with the protocol from \cite{ruiz2023dreambooth}, we carry out a quantitative analysis that examines subject fidelity (DINO~\cite{caron2021emerging}, CLIP-I~\cite{radford2021learning}), adherence to text prompts (CLIP-T~\cite{radford2021learning}), and image diversity (LPIPS~\cite{zhang2018unreasonable}). For CLIP-I, we compute the mean pairwise cosine similarity between CLIP embeddings from generated and real images. The DINO metric is similar but uses embeddings from ViT S/16 DINO. The CLIP-T measure involves calculating the mean cosine similarity between embeddings from the text prompt and those from generated images. Additionally, we analyze the average LPIPS cosine similarity across images generated for the same subject using an identical text prompt. Table~\ref{table:dreambooth_final} demonstrates that both COFT and OFT achieve significantly better results than DreamBooth and LoRA in terms of DINO and CLIP-I, and perform as well as or marginally better in prompt fidelity and diversity scores. To mitigate the effect of randomness, each method is evaluated by finetuning the same pretrained model 30 times, each with a unique random seed. Altogether, the findings indicate that the OFT framework not only improves on convergence and stability but also secures consistently stronger final results.

\textbf{Qualitative comparison}. To provide clearer insight into the advantages of OFT, Figure~\ref{fig:dreambooth_final} displays several randomly selected examples showcasing subject-driven generation. To ensure an equitable comparison, all outputs are produced using the same finetuned model for each approach, without optimizing individual text prompts for the final results. For every method, the model exhibiting the highest validation CLIP scores is chosen. As shown in Figure~\ref{fig:dreambooth_final}, both OFT and COFT are able to robustly maintain semantic fidelity to the original subject, whereas LoRA frequently struggles to retain subject identity (\emph{e.g.}, in the bowl example, LoRA completely fails to preserve the subject). Additionally, both OFT and COFT demonstrate superior text prompt controllability, while DreamBooth, though proficient at maintaining subject identity, frequently does not adhere to the prompt in generated images (\emph{e.g.}, see the bowl example in the first row). This qualitative analysis confirms that OFT provides enhanced control and subject preservation simultaneously. Furthermore, OFT's effectiveness is largely invariant to the number of iterations, maintaining strong performance even when the iteration count is increased—a trait not shared by DreamBooth or LoRA. Additional qualitative samples are provided in Appendix~\ref{app:more_qualitative}. We also include a human evaluation in Appendix~\ref{app:human}, which further corroborates the superiority of OFT.

\subsection{Controllable Generation}

\textbf{Settings}. As baselines, we incorporate ControlNet~\cite{zhang2023adding}, T2I-Adapter~\cite{mou2023t2i}, and LoRA~\cite{hulora2022}. The main paper evaluates performance on three demanding controllable generation tasks: Canny edge to image~(C2I) using the COCO dataset~\cite{lin2014microsoft}, segmentation map to image~(S2I) using the ADE20K dataset~\cite{zhou2017scene}, and landmark to face~(L2F) on the CelebA-HQ dataset~\cite{karras2018progressive,xia2021tedigan}. Each approach is employed to finetune Stable Diffusion~(SD) v1.5 for 20 epochs across these datasets. Further experimental results are presented in Appendix~\ref{app:more_qualitative}, \ref{app:more_control_task}, and \ref{app:failure}.

\textbf{Convergence}. We assess the rate at which ControlNet, T2I-Adapter, LoRA, and COFT converge on the L2F task using both quantitative and qualitative methods. For quantitative analysis, we utilize the mean $\ell_2$ distance between predicted face landmarks and control face landmarks as our metric. Figure~\ref{fig:face_convergence} displays how the face landmark error decreases as each model is finetuned over varying epochs. It is evident that COFT and OFT reach lower error rates significantly quicker than competing approaches. For instance, LoRA only attains OFT’s 8-th epoch performance after undergoing 20 epochs of training. Notably, OFT and COFT maintain a comparable number of trainable parameters to LoRA—substantially less than ControlNet—yet attain improved convergence efficiency relative to existing baselines. Additionally, Figure~\ref{fig:teaser} further demonstrates OFT’s rapid convergence: the rightmost example illustrates that OFT effectively learns with merely 5\% of the ADE20K dataset, underscoring its data efficiency when compared to both ControlNet and LoRA. For our qualitative comparisons, we primarily highlight OFT, COFT, and ControlNet, since ControlNet’s final landmark error is closest to that of our models. As depicted in Figure~\ref{fig:face_convergence_qual}, both OFT and COFT show steady convergence; generated face poses align more closely with control landmarks over time. In contrast, ControlNet exhibits an abrupt emergence of controllability post the 8-th epoch, which is consistent with the sudden convergence phenomenon reported in \cite{zhang2023adding}. This consistent and gradual convergence characteristic of OFT simplifies practical deployment, as its training process remains stable and more predictable, making the tuning of OFT’s hyperparameters a more straightforward task.

\textbf{Quantitative comparison}. To assess the effectiveness of controllable generation, we propose a control consistency metric that measures alignment between the control signal extracted from the synthesized image and the corresponding original input control signal. Specifically, for the C2I task, the metric uses Intersection over Union (IoU) and F1 score, while for the S2I task, mean IoU along with mean and overall accuracy are calculated. For the L2F task, the evaluation is based on the average $\ell_2$ distance between predicted and control landmarks. A comprehensive description of these metrics is provided in Appendix~\ref{app:settings}. All methods under comparison are evaluated with their optimal hyperparameter configurations. As summarized in Table~\ref{table:control_gen}, OFT and COFT both demonstrate superior control accuracy relative to alternative approaches. In contrast, adapter-based methods such as T2I-Adapter and ControlNet display slower convergence and underperform in terms of final control quality. Notably, LoRA surpasses ControlNet on the S2I task but lags behind in C2I and L2F tasks; overall, LoRA’s maximal performance is modest even with meticulous hyperparameter optimization. By comparison, our empirical findings indicate that the OFT framework continues to benefit from increased parameterization, suggesting its performance has yet to plateau. Importantly, the quantitative evaluation protocol we introduce for controllable generation constitutes a novel contribution, as it provides a robust and precise assessment of fine-tuned model control (subject to the accuracy limits of the deployed segmentation or detection backbone).

\textbf{Qualitative comparison}. We further present a qualitative assessment of OFT and COFT alongside leading approaches such as ControlNet, T2I-Adapter, and LoRA. As depicted in Figure~\ref{fig:control_final}, random image samples reveal that OFT and COFT not only produce visually convincing and high-quality results, but also afford precise controllability. Within the S2I task, LoRA is unable to adhere to the input segmentation map, failing to generate compliant outputs, whereas ControlNet, OFT, and COFT successfully guide image synthesis according to the segmentation. Notably, OFT and COFT are capable of generating superior image fidelity and richer, more realistic details with improved structural coherence compared to ControlNet, despite requiring significantly fewer model parameters. On the C2I task, OFT and COFT effectively synthesize semantically coherent images from coarse Canny edge maps, outperforming both T2I-Adapter and LoRA. For the L2F task, our method demonstrates the highest accuracy in pose guidance for synthesized faces, even in cases involving complex facial orientations. Across all three control benchmarks, OFT and COFT consistently deliver better qualitative results relative to state-of-the-art baselines, underscoring the efficacy of the OFT framework for controlled image generation. To broaden this qualitative analysis, we include additional results covering all three tasks in Appendix~\ref{app:control_exp}, and further illustrate OFT’s versatility on extra control scenarios—such as dense pose to human body, sketch to image, and depth to image—in Appendix~\ref{app:more_control_task}.

\section{Concluding Remarks and Open Problems}

Inspired by the critical role of angular relationships between neurons in shaping visual semantics, we introduce orthogonal finetuning (OFT), a straightforward yet powerful approach to refining text-to-image diffusion models. Our method is specifically applied to two prominent tasks: subject-driven generation and controllable generation. Relative to prior approaches, OFT achieves enhanced control and greater finetuning stability while requiring fewer parameters. Notably, OFT does not incur any additional cost during inference, thereby offering a model that can be efficiently deployed in practice.

OFT also brings forth several notable open questions. To ensure orthogonality, OFT utilizes Cayley parametrization, which entails a matrix inversion. This requirement can restrict the scalability of OFT; while block diagonal parametrization provides a partial remedy, finding a differentiable and faster method for this matrix inversion remains unresolved. Moreover, OFT presents intriguing possibilities for compositionality: the orthogonal matrices learned through different OFT finetuning tasks can be multiplied together and the result remains orthogonal. A compelling research direction is to investigate whether this collection of orthogonal matrices retains the specialized knowledge from each task. Lastly, OFT's parameter efficiency is mainly dictated by the block diagonal design, which can introduce bias and constrain flexibility. Thus, devising strategies to further improve parameter efficiency while reducing bias remains a key challenge for future work.

\end{document}